% =============================================================================
% UQFF v5.78  Section Template :: T-PRED
% Open-Prediction papers (predicted-but-not-yet-observed; 9999 sentinels)
% -----------------------------------------------------------------------------
% USAGE
%   1. Copy to whitepapers/PAPER_<NNNN>_<slug>.tex
%   2. Replace every <PLACEHOLDER>
%   3. The closure tracker tags err==9999 rows from this template into the
%      OPEN_PREDICTIONS ledger tab, NOT into the high-error physics bucket.
% =============================================================================
\documentclass[11pt]{article}
\usepackage[margin=1in]{geometry}
\usepackage{amsmath,amssymb,amsthm}
\usepackage{booktabs,longtable,array}
\usepackage{listings,xcolor}
\usepackage[colorlinks=true,linkcolor=blue,urlcolor=blue,citecolor=blue]{hyperref}
\lstset{basicstyle=\ttfamily\small,breaklines=true,frame=single,
        language=Python,showstringspaces=false,columns=fullflexible}

\title{<PAPER TITLE>: \\ Open UQFF Prediction (Pre-Observation), v5.78}
\author{Star-Magic UQFF \\ PAPER\_<NNNN> \,/\, Session <SESSION>}
\date{<DATE>}

% NOTE: predicted is a finite number; observed=9999 is the OPEN-PREDICTION sentinel.
% Error %=9999 is also a sentinel and MUST NOT be aggregated as a physics gap.
% CLOSURE :: <closure_label> :: predicted=<X> observed=9999 error_pct=9999
% TEMPLATE :: T-PRED
% CANONICAL :: open_prediction = true

\begin{document}
\maketitle

\begin{abstract}
The UQFF v5.78 framework predicts the value of <observable>; no published
measurement currently exists at the required precision. This paper
documents the prediction, the derivation chain, the observational
program(s) that could test it, and the threshold at which the prediction
falsifies.
\end{abstract}

\section{Prediction}
\begin{equation}
  \boxed{\;<UQFF observable> \;=\; <numerical prediction with units>\;}
\end{equation}
Derivation chain: <PAPER\_NN $\to$ PAPER\_MM $\to$ this paper>.

\section{Falsification Criterion}
The prediction is falsified at the $5\sigma$ level if a future measurement
yields a value outside $<P>\pm 5\sigma_P$ where $\sigma_P$ is the propagated
theoretical uncertainty from the canonical constants:
\begin{itemize}
  \item $\rho_{\mathrm{vac,SCm}}$: structural (zero theoretical uncertainty).
  \item $\beta_i \approx 0.603$ ($\pm 0.0004\%$ from the closure hunt).
  \item $\kappa = 5\times 10^{-4}/\mathrm{day}$ (exact rational $5/10000$).
  \item $[SSq] = 0.57$ (exact rational $57/100$).
\end{itemize}
Propagated $\sigma_P$ = <value>.

\section{Suggested Observational Tests}
\begin{longtable}{lll}
\toprule
Probe / Instrument & Target sensitivity & Earliest data release \\
\midrule
<e.g. LISA, CTA, SKA-low, Roman, ELT, LIGO O5, future LHCb run> & <value> & <year> \\
\bottomrule
\end{longtable}

\section{Status in the Ledger}
This closure carries \texttt{observed=9999, error\_pct=9999} as a sentinel
and is filed in the \texttt{OPEN\_PREDICTIONS} tab of the ledger triage.
Do \emph{not} aggregate this into the high-error physics bucket.

\section{Cross-references}
\begin{itemize}
  \item Sister open predictions: <e.g. $A_s$ at S764, $\eta_b$ at S764>.
  \item Canonical source: \texttt{dpm\_vacuum\_manifold.py} v3.0.
\end{itemize}

\bibliographystyle{plain}
\begin{thebibliography}{9}
\bibitem{ref1} <DOI/arXiv for the instrument design report or proposal>.
\end{thebibliography}

\end{document}

